\chapter{Perbandingan Senilai, Kecepatan dan Rata-rata}\label{ch:g8:speed}

\omterm{def:g7:prop:table}{Perbandingan senilai} (\cref{ch:g7:prop}) kini bertemu dunia fisik:
\omterm{def:g8:speed:def}{kecepatan}, debit, pengubahan satuan --- besaran yang berupa
\emph{\omterm{def:g3:division:remainder}{hasil bagi}} dua besaran lain. Bab ini ditutup dengan \omterm{def:g8:speed:average}{rata-rata
berbobot}, tempat penalaran yang \omterm{def:g6:propdata:def}{sebanding} mencegah kekeliruan klasik.

\section{Kecepatan rata-rata}

\begin{definition}[Kecepatan rata-rata]\label{def:g8:speed:def}
\emph{Kecepatan rata-rata}\index{kecepatan} sebuah perjalanan adalah \omterm{def:g3:division:remainder}{hasil bagi}
\[
v = \frac{d}{t}
\qquad
\left(\text{jarak dibagi lamanya}\right),
\]
dalam kilometer per jam (km/jam), atau meter per detik (m/detik),
\dots\ Rumusnya dapat dibalik: $d = v \times t$ dan $t = \frac{d}{v}$.
\end{definition}

\begin{example}\label{ex:g8:speed:basic}
Sebuah kereta menempuh $270$ km dalam $1$ jam $30$ menit. Ubahlah
waktunya lebih dahulu: $1$ jam $30$ menit $= 1.5$ jam. Lalu
\[
v = \frac{270}{1.5} = 180 \text{ km/jam}.
\]
Jarak pada $90$ km/jam selama $2$ jam $20$ menit $= \frac73$ jam:
$d = 90 \times \frac73 = 210$ km. Waktu untuk $35$ km pada
$14$ km/jam: $t = \frac{35}{14} = 2.5$ jam $= 2$ jam $30$ menit.
\end{example}

\begin{remark}[Menit bukan bilangan desimal]
$1$ jam $30$ menit adalah $1.5$ jam, tetapi $1$ jam $20$ menit
\emph{bukan} $1.2$ jam: melainkan $1 + \frac{20}{60} = \frac43 \approx
1.33$ jam. Ubahlah menitnya menjadi \omterm{def:g6:fractions:def}{pecahan} dari satu jam
($60$ menit $= 1$ jam) sebelum membagi.
\end{remark}

\begin{omfigure}
\begin{tikzpicture}
\begin{axis}[omaxis,
  width=8.4cm, height=5.6cm,
  xmin=0, xmax=3.4, ymin=0, ymax=270,
  xlabel={waktu (jam)}, ylabel={jarak (km)},
  xtick={1,2,3}, ytick={80,160,240}]
  \addplot[omDef, thick, domain=0:3.2] {80*x};
  \addplot[omThm, thick, domain=0:3.2] {50*x};
  \node[omDef, font=\small, rotate=36] at (axis cs:2.35,215)
    {$80$ km/jam};
  \node[omThm, font=\small, rotate=25] at (axis cs:2.6,155)
    {$50$ km/jam};
\end{axis}
\end{tikzpicture}

{\small Pada \omterm{def:g8:speed:def}{kecepatan} tetap, jarak \omterm{def:g6:propdata:def}{sebanding} dengan waktu: yaitu
garis lurus melalui titik asalnya, yang kecuramannya \emph{adalah}
\omterm{def:g8:speed:def}{kecepatannya}.}
\end{omfigure}

\begin{method}[Mengubah kecepatan]\label{met:g8:speed:convert}
Untuk mengubah km/jam menjadi m/detik: $1$ km $= 1000$ m dan
$1$ jam $= 3600$ detik, jadi
\[
v \text{ km/jam} = \frac{v \times 1000}{3600} \text{ m/detik}
= \frac{v}{3.6} \text{ m/detik}.
\]
Untuk mengubah m/detik menjadi km/jam, kalikan dengan $3.6$.
\end{method}

\begin{example}\label{ex:g8:speed:convert}
$90$ km/jam $= \frac{90}{3.6} = 25$ m/detik. Pelari cepat yang
menempuh $100$ m dalam $10$ detik bergerak pada $10$ m/detik
$= 36$ km/jam.
\end{example}

\section{Besaran yang berupa hasil bagi}

\begin{example}[Debit, harga per kilo]\label{ex:g8:speed:quotients}
\omterm{def:g8:speed:def}{Kecepatan} punya banyak sepupu, yang semuanya diperlakukan dengan cara
yang sama:
\begin{itemize}
  \item sebuah keran mengisi $48$ L dalam $4$ menit: debitnya
        $\frac{48}{4} = 12$ L/menit, jadi mengisi bak $150$ L makan
        waktu $\frac{150}{12} = 12.5$ menit;
  \item $0.6$ kg keju berharga $9$ euro: harganya
        $\frac{9}{0.6} = 15$ euro tiap kg;
  \item sebuah mobil memakai $6.3$ L untuk $90$ km: konsumsinya
        $\frac{6.3}{90} \times 100 = 7$ L tiap $100$ km.
\end{itemize}
Kenalilah \omterm{def:g3:division:remainder}{hasil baginya}, lalu rumus tiga arahnya
($q = \frac ab$, $a = qb$, $b = \frac aq$) mengerjakan sisanya.
\end{example}

\section{Rata-rata berbobot}

\begin{definition}[Rata-rata berbobot]\label{def:g8:speed:average}
Ketika nilai $v_1, v_2, \dots$ datang bersama \omterm{def:g7:stats:frequency}{frekuensi} (atau bobot)
$n_1, n_2, \dots$, \emph{rata-rata berbobotnya}\index{rata-rata berbobot}
adalah
\[
\bar v = \frac{n_1 v_1 + n_2 v_2 + \dots}{n_1 + n_2 + \dots} :
\]
\omterm{def:g1:addition:def}{jumlah} nilainya, dibagi \omterm{def:g1:addition:def}{jumlah} bobotnya.
\end{definition}

\begin{example}\label{ex:g8:speed:average}
Sebuah ujian dikerjakan dua kelompok: kelompok 1, $12$ murid,
rata-ratanya $11$; kelompok 2, $18$ murid, rata-ratanya $14$.
Rata-rata seluruh kelasnya:
\[
\bar v = \frac{12 \times 11 + 18 \times 14}{12 + 18}
= \frac{132 + 252}{30} = \frac{384}{30} = 12.8 .
\]
\emph{Bukan} $\frac{11 + 14}{2} = 12.5$: kelompok yang lebih besar
menarik rata-ratanya ke arahnya sendiri --- bobotnya berpengaruh.
\end{example}

\begin{example}[Kecepatan rata-rata pada dua ruas]\label{ex:g8:speed:twolegs}
Seorang pesepeda menempuh $30$ km pada $30$ km/jam, lalu $30$ km pada
$15$ km/jam. \omterm{def:g8:speed:def}{Kecepatan rata-rata} seluruh perjalanannya \emph{bukan}
$\frac{30 + 15}{2} = 22.5$ km/jam. Hitunglah dengan definisinya:
\begin{enumerate}
  \item waktunya: $\frac{30}{30} = 1$ jam, lalu $\frac{30}{15} = 2$
        jam; seluruhnya $3$ jam;
  \item jarak seluruhnya $60$ km, jadi
        $v = \frac{60}{3} = 20$ km/jam.
\end{enumerate}
Ruas yang lambat berlangsung lebih lama, jadi bobotnya lebih besar ---
rata-rata \omterm{def:g8:speed:def}{kecepatan} harus dibobot menurut \emph{waktu}.
\end{example}

\section{Latihan}

\begin{exercise}[$\star$]\label{exo:g8:speed:1}
Hitunglah \omterm{def:g8:speed:def}{kecepatan} rata-ratanya: $150$ km dalam $2$ jam; $27$ km
dalam $45$ menit; $100$ m dalam $12.5$ detik.
\end{exercise}

\begin{exercise}[$\star$]\label{exo:g8:speed:2}
Hitunglah jaraknya: $2$ jam pada $85$ km/jam; $40$ menit pada
$90$ km/jam. Hitunglah waktunya: $315$ km pada $90$ km/jam (dalam jam
dan menit).
\end{exercise}

\begin{exercise}[$\star$]\label{exo:g8:speed:3}
Ubahlah: $72$ km/jam menjadi m/detik; $5$ m/detik menjadi km/jam;
$108$ km/jam menjadi m/detik.
\end{exercise}

\begin{exercise}[$\star$]\label{exo:g8:speed:4}
Sebuah keran mengalirkan $15$ L/menit. Berapa lama untuk mengisi
tangki $600$ L? Berapa liter dalam $2$ jam $30$?
\end{exercise}

\begin{exercise}[$\star$]\label{exo:g8:speed:5}
Mana yang lebih menguntungkan: $1.2$ kg apel seharga $3.30$ euro, atau
$0.8$ kg seharga $2.16$ euro? Bandingkan harganya tiap kilogram.
\end{exercise}

\begin{exercise}[$\star$]\label{exo:g8:speed:6}
Sebuah kelas berisi $25$ murid punya rata-rata $12.4$ pada sebuah
ujian; kelas lain berisi $35$ murid punya rata-rata $10$. Hitunglah
rata-rata kedua kelas itu bersama-sama.
\end{exercise}

\begin{exercise}[$\star$]\label{exo:g8:speed:7}
Bunyi merambat kira-kira $340$ m/detik. Kamu melihat kilatan petir
lalu mendengar guruhnya $6$ detik kemudian. Berapa jauh petirnya
menyambar (sampai $100$ m terdekat)?
\end{exercise}

\begin{exercise}[$\star\star$]\label{exo:g8:speed:8}
Seorang pendaki berjalan $2$ jam pada $5$ km/jam, beristirahat
$30$ menit, lalu berjalan $1$ jam $30$
pada $4$ km/jam. Hitunglah jarak seluruhnya dan \omterm{def:g8:speed:def}{kecepatan} rata-ratanya
\emph{termasuk istirahatnya} (jarak seluruhnya dibagi waktu yang berlalu).
\end{exercise}

\begin{exercise}[$\star\star$]\label{exo:g8:speed:9}
Nilai dengan koefisien: seorang murid memperoleh $15$ (koefisien $3$),
$9$ (koefisien $2$) dan $12$ (koefisien $1$). Hitunglah \omterm{def:g8:speed:average}{rata-rata
berbobotnya}. Rata-rata biasa (tanpa bobot) yang mana yang akan
diperolehnya, dan mengapa keduanya berbeda?
\end{exercise}

\begin{exercise}[$\star\star$]\label{exo:g8:speed:10}
Sebuah mobil menempuh $120$ km pada $80$ km/jam lalu $120$ km pada
$120$ km/jam.
\begin{enumerate}
  \item Hitunglah lamanya tiap ruasnya, lalu \omterm{def:g8:speed:def}{kecepatan} rata-ratanya
        pada seluruh perjalanannya.
  \item Jelaskan mengapa jawabannya kurang dari $100$ km/jam, yaitu
        titik tengah kedua \omterm{def:g8:speed:def}{kecepatannya}.
\end{enumerate}
\end{exercise}

\begin{exercise}[$\star\star\star$]\label{exo:g8:speed:11}
Dua desa berjarak $18$ km. Ana berangkat dari desa pertama pukul
$10{:}00$ dengan berjalan $5$ km/jam menuju desa kedua; Bayu berangkat
dari desa kedua pada saat yang sama dengan berjalan $4$ km/jam menuju
desa pertama. Pukul berapa mereka bertemu, dan berapa jaraknya dari
desa Ana? (Bersama-sama mereka menutup celahnya pada $5 + 4$ km/jam.)
\end{exercise}

\section{Soal: Rata-rata harmonik, atau mengapa perjalanan pulang
merusak rata-ratanya}

\begin{problem}[{Soal akhir pekan --- \omterm{def:g8:speed:def}{kecepatan} rata-rata pada jarak
yang sama adalah $\frac{2 v_1 v_2}{v_1 + v_2}$, dan tidak pernah
mengalahkan titik tengah kedua \omterm{def:g8:speed:def}{kecepatannya}}]
\label{pb:g8:speed:1}
\Cref{ex:g8:speed:twolegs} dan \cref{exo:g8:speed:10} keduanya
berakhir dengan kejutan yang sama: berangkat pada satu \omterm{def:g8:speed:def}{kecepatan},
kembali pada \omterm{def:g8:speed:def}{kecepatan} yang lain, dan \omterm{def:g8:speed:def}{kecepatan} rata-ratanya mendarat
\emph{di bawah} titik tengah keduanya. Soal ini mencari rumus persis
di balik kejutannya --- yaitu jenis rata-rata yang baru,
\emph{rata-rata harmonik}\index{rata-rata harmonik} --- membuktikan
bahwa kejutannya adalah sebuah teorema, lalu mendorongnya ke
ujung logisnya: pengejaran yang tidak dapat dicapai \omterm{def:g8:speed:def}{kecepatan} mana pun
di dunia.

\medskip
\noindent\textbf{Bagian I --- Rumus untuk perjalanan pulang pergi.}
Sebuah perjalanan menempuh jarak $d$ pada \omterm{def:g8:speed:def}{kecepatan} $v_1$, lalu jarak
$d$ yang sama lagi pada \omterm{def:g8:speed:def}{kecepatan} $v_2$ (semuanya bilangan positif).
\begin{enumerate}
  \item Nyatakan kedua lamanya, lalu lama seluruhnya, sebagai
        \omterm{def:g6:fractions:def}{pecahan} yang memuat $d$, $v_1$, $v_2$.
  \item \omterm{def:g8:speed:def}{Kecepatan rata-rata} seluruh perjalanannya adalah
        $\bar v = \dfrac{2d}{\ \dfrac{d}{v_1} + \dfrac{d}{v_2}\ }$
        --- yaitu \omterm{def:g6:fractions:def}{pecahan} bertingkat
        (\cref{ex:g8:fractions:complex}). Sederhanakan lalu buktikan:
        \[
        \bar v = \frac{2\, v_1 v_2}{v_1 + v_2} .
        \]
  \item Periksalah rumusnya terhadap kedua perhitungan yang dikutip
        di atas: $v_1 = 30$, $v_2 = 15$ km/jam
        (\cref{ex:g8:speed:twolegs}), dan $v_1 = 80$,
        $v_2 = 120$ km/jam (\cref{exo:g8:speed:10}).
  \item Jarak $d$ sudah lenyap dari rumusnya. Apa artinya itu,
        secara nyata, bagi pesepedanya?
  \item Sekarang andaikan perjalanannya justru menghabiskan
        \emph{waktu} $T$ yang sama pada tiap \omterm{def:g8:speed:def}{kecepatannya}. Tunjukkan
        bahwa \omterm{def:g8:speed:def}{kecepatan} rata-ratanya lalu menjadi
        $\frac{v_1 + v_2}{2}$, yaitu titik tengah yang biasa.
        Dalam bahasa \cref{def:g8:speed:average}: dengan besaran
        apa rata-rata \omterm{def:g8:speed:def}{kecepatan} harus selalu dibobot, dan apa yang
        berubah antara kedua keadaannya?
\end{enumerate}

\medskip
\noindent\textbf{Bagian II --- Harmonik melawan aritmetika.}
Untuk dua bilangan positif $v_1$, $v_2$, tetapkan
\[
H = \frac{2\, v_1 v_2}{v_1 + v_2}
\quad\text{(\emph{rata-rata harmonik}),}
\qquad
A = \frac{v_1 + v_2}{2}
\quad\text{(\emph{rata-rata aritmetika}).}
\]
\begin{enumerate}[resume]
  \item Hitunglah $H$ dan $A$ untuk pasangan $(30, 15)$ dan
        $(80, 120)$. Rata-rata yang mana dari keduanya yang menang
        setiap kali?
  \item Taruhlah $A - H$ pada \omterm{def:g4:fractions:def}{penyebut} bersama $2(v_1 + v_2)$,
        jabarkan \omterm{def:g4:fractions:def}{pembilangnya} dengan kesamaan istimewa
        (\cref{pb:g8:equations:1}), lalu buktikan:
        \[
        A - H = \frac{(v_1 - v_2)^2}{2\,(v_1 + v_2)} .
        \]
  \item Simpulkan teorema di balik semua kejutannya: untuk \omterm{def:g8:speed:def}{kecepatan}
        yang positif, $H \leq A$ selalu, dengan kesamaannya persis
        ketika $v_1 = v_2$. Mengapa kuadrat pada \omterm{def:g4:fractions:def}{pembilangnya}
        menyelesaikan perkaranya?
  \item Buktikan kesamaan yang anggun $H \times A = v_1 \times v_2$:
        kedua rata-ratanya mengalikan kembali menjadi \omterm{def:g2:mult:def}{hasil kali}
        aslinya.
  \item Sebuah mobil menempuh $60$ km pada $40$ km/jam dan $60$ km
        pada $60$ km/jam. Ramalkan \omterm{def:g8:speed:def}{kecepatan} rata-ratanya dengan
        rumusnya, lalu benarkan lewat jalan panjangnya, dengan waktu
        dan jarak seluruhnya.
\end{enumerate}

\medskip
\noindent\textbf{Bagian III --- Pengejaran yang mustahil.}
\begin{enumerate}[resume]
  \item Seorang pesepeda mengikuti lomba $30$ km dengan harapan
        rata-ratanya $30$ km/jam. Paling lama berapa waktu seluruhnya
        yang boleh dipakai lombanya? Ia menempuh $15$ km pertamanya
        pada $15$ km/jam: berapa sisa jatah waktunya?
  \item Tunjukkan bahwa rata-rata yang diharapkannya sudah menjadi
        mustahil secara harfiah: hitunglah \omterm{def:g8:speed:def}{kecepatan} rata-ratanya
        kalau ia menempuh separuh keduanya pada $90$ km/jam, lalu
        jelaskan mengapa \omterm{def:g8:speed:def}{kecepatan} yang raksasa sekalipun tetap
        membuatnya kurang dari
        $30$ km/jam.
  \item Ia menurunkan sasarannya menjadi $20$ km/jam. \omterm{def:g8:speed:def}{Kecepatan} apa
        pada $15$ km yang kedua yang mencapainya? Periksalah
        jawabanmu dengan rumus \omterm{pb:g8:speed:1}{rata-rata harmoniknya}.
  \item Keran dingin sendirian mengisi sebuah bak dalam $20$ menit;
        keran panas sendirian mengisinya dalam $30$ menit. Berapa
        bagian baknya yang diisi kedua kerannya bersama-sama tiap
        menitnya, dan berapa lama waktu yang dipakainya?
        Bandingkan jawabanmu dengan $H(20, 30)$ --- apa yang kamu
        perhatikan?
  \item Sebuah pesawat menerbangi rute yang sama pergi dan pulang:
        $900$ km/jam bersama angin, $700$ km/jam melawannya.
        Hitunglah \omterm{def:g8:speed:def}{kecepatan} rata-rata pulang perginya, lalu jelaskan
        dalam satu kalimat mengapa angin, ke mana pun ia bertiup,
        selalu \emph{memperpanjang} perjalanan pulang pergi.

\end{enumerate}
\end{problem}
